% AgroVoltaic --- formulario. Todas las formulas que usa el sistema, con el
% archivo donde vive cada una.
% Compilar: xelatex -interaction=nonstopmode formulas.tex (dos veces, por el indice).
\documentclass[11pt,a4paper]{article}
% Preambulo compartido: paquetes, paleta, estilos de TikZ y macros.
% Responsabilidad unica: configurar el documento. No contiene contenido.

% Tipografia (requiere xelatex; ver README). Pagella para el cuerpo, Heros para
% titulos y diagramas, JetBrains Mono para codigo. Las tres estan instaladas:
% las dos primeras vienen con texlive-fontsrecommended, la tercera es del sistema.
\usepackage{fontspec}
\usepackage{unicode-math}
% Las TeX Gyre viven en texmf y NO estan registradas en fontconfig: se cargan por
% nombre de archivo, no por nombre de familia.
\setmainfont{texgyrepagella}[
  Extension=.otf, UprightFont=*-regular, ItalicFont=*-italic,
  BoldFont=*-bold, BoldItalicFont=*-bolditalic]
\setsansfont{texgyreheros}[
  Extension=.otf, UprightFont=*-regular, ItalicFont=*-italic,
  BoldFont=*-bold, BoldItalicFont=*-bolditalic, Scale=MatchLowercase]
\setmonofont{JetBrains Mono}[Scale=MatchLowercase]
\setmathfont{texgyrepagella-math.otf}
\usepackage[spanish,es-noquoting,es-tabla]{babel}
% expansion no existe en xetex; protrusion si, y es la que mas se nota.
\usepackage[protrusion=true,expansion=false]{microtype}
\usepackage{geometry}
\geometry{a4paper,margin=2.4cm,headsep=12pt}
\usepackage{amsmath}
\usepackage{booktabs}
\usepackage{longtable}
\usepackage{array}
\usepackage{enumitem}
\usepackage{tcolorbox}
\usepackage{tikz}
\usepackage{fancyhdr}
\usepackage{titlesec}
\usepackage[hidelinks]{hyperref}

\usetikzlibrary{positioning,shapes.geometric,arrows.meta,calc,fit,backgrounds}
\tcbuselibrary{skins,breakable}

% --- Paleta -----------------------------------------------------------------
\definecolor{solar}{HTML}{B26B00}   % naranja apagado: fuentes fisicas / energia
\definecolor{tierra}{HTML}{4F6B3A}  % verde: agro / AgroDash
\definecolor{agua}{HTML}{20566E}    % azul: almacenamiento
\definecolor{tinta}{HTML}{1F2328}   % texto
\definecolor{humo}{HTML}{6B7280}    % texto secundario
\definecolor{papel}{HTML}{F5F3EE}   % fondo de cajas
\definecolor{borde}{HTML}{D6D2C8}

\hypersetup{colorlinks=false,pdftitle={AgroVoltaic --- Arquitectura},
            pdfauthor={Proyecto AgroVoltaic (TEC)}}

% --- Titulos ----------------------------------------------------------------
\titleformat{\section}{\Large\bfseries\color{agua}}{\thesection}{0.7em}{}
\titleformat{\subsection}{\large\bfseries\color{tinta}}{\thesubsection}{0.7em}{}
\titlespacing*{\section}{0pt}{2.2ex plus 1ex minus .2ex}{1.2ex plus .2ex}

\pagestyle{fancy}
\fancyhf{}
\fancyhead[L]{\small\color{humo}AgroVoltaic --- Arquitectura del sistema}
\fancyhead[R]{\small\color{humo}\thepage}
\renewcommand{\headrulewidth}{0.4pt}

% --- Macros -----------------------------------------------------------------
\newcommand{\cd}[1]{\texttt{\small #1}}          % identificador / ruta / comando
% \nota: texto secundario. OJO: no admite \\ adentro — el salto de linea de un
% nodo de TikZ es una celda de alineacion y no cruza el grupo. Una \nota por linea.
\newcommand{\nota}[1]{{\color{humo}\small #1}}

\newtcolorbox{clave}[1][]{
  enhanced, breakable, colback=papel, colframe=agua, boxrule=0.6pt,
  left=8pt, right=8pt, top=6pt, bottom=6pt, arc=2pt,
  fonttitle=\bfseries\small, coltitle=white, colbacktitle=agua, #1}

\newtcolorbox{aviso}[1][]{
  enhanced, breakable, colback=white, colframe=solar, boxrule=0.6pt,
  left=8pt, right=8pt, top=6pt, bottom=6pt, arc=2pt,
  fonttitle=\bfseries\small, coltitle=white, colbacktitle=solar, #1}

% --- Estilos de diagrama ----------------------------------------------------
\tikzset{
  comp/.style={
    draw=agua, fill=papel, line width=0.7pt, rounded corners=3pt,
    align=center, inner sep=6pt, minimum height=10mm, font=\small},
  compx/.style={comp, draw=humo, fill=white},          % componente externo
  fuente/.style={comp, draw=solar, fill=solar!8},      % fuente fisica / dato crudo
  agro/.style={comp, draw=tierra, fill=tierra!8},      % region Cartago
  almacen/.style={
    draw=agua, fill=agua!8, line width=0.7pt, cylinder, shape border rotate=90,
    aspect=0.16, align=center, inner sep=6pt, minimum height=13mm, font=\small},
  etapa/.style={comp, minimum width=27mm},
  flecha/.style={-{Stealth[length=2.4mm]}, line width=0.7pt, draw=tinta},
  flechad/.style={flecha, dashed, draw=humo},
  etiq/.style={font=\scriptsize\itshape, color=humo, inner sep=2pt, fill=white},
  grupo/.style={draw=borde, line width=0.6pt, rounded corners=4pt, inner sep=8pt},
  % Sin guiones dentro de los diagramas: partir "ana-lizador" en una caja de
  % 27 mm se lee peor que dejar la caja un poco mas holgada.
  every picture/.append style={execute at begin picture={%
    \hyphenpenalty=10000\relax\exhyphenpenalty=10000\relax}},
  titgrupo/.style={font=\scriptsize\bfseries, color=humo},
}


% El preambulo es compartido con el documento de arquitectura: se le cambia solo
% el encabezado, para que cada pagina diga de que documento es.
\fancyhead[L]{\small\color{humo}AgroVoltaic $\cdot$ Formulario}
\hypersetup{pdftitle={AgroVoltaic. Formulario}}

\begin{document}

\begin{titlepage}
\centering
\vspace*{3.5cm}
{\color{humo}\large PROYECTO AGROVOLTAICO $\cdot$ TEC COSTA RICA\par}
\vspace{1.2cm}
{\Huge\bfseries\color{agua} AgroVoltaic\par}
\vspace{0.4cm}
{\LARGE Formulario\par}
\vspace{1.2cm}
{\large\color{tinta} Todo el cálculo del sistema, en un solo lugar\par}
\vfill
\begin{minipage}{0.82\textwidth}
\small\color{tinta}
Cada fórmula que el sistema aplica a un número, con el símbolo que usa, la condición
bajo la que vale y el archivo donde está implementada. No hay ninguna que no esté en
el código, ni ninguna en el código que no esté acá.
\end{minipage}
\vfill
{\color{humo} \today\par}
\end{titlepage}

\tableofcontents
\newpage

% =============================================================================
\section{Símbolos}

\begin{center}
\small
\begin{tabular}{@{}p{26mm}p{88mm}p{22mm}@{}}
\toprule
\textbf{Símbolo} & \textbf{Qué es} & \textbf{Unidad} \\
\midrule
$\mathrm{GHI}$ & Irradiancia global sobre plano horizontal, la que mide el piranómetro & W/m$^2$ \\
$\mathrm{GHI_{cs}}$ & La misma, pero teórica, con cielo perfectamente despejado & W/m$^2$ \\
$\mathrm{kt^*}$ & Índice de claridad: qué fracción de la luz posible llegó de verdad & adimensional \\
$\mathrm{DNI}$ & Componente directa, la que viene del disco solar & W/m$^2$ \\
$\mathrm{DHI}$ & Componente difusa, la que llega del resto del cielo & W/m$^2$ \\
$\mathrm{POA}$ & Irradiancia sobre el plano del arreglo (\emph{plane of array}) & W/m$^2$ \\
$\rho$ & Albedo: fracción de luz que el suelo refleja & adimensional \\
$\varphi$ & Factor de bifacialidad: cuánto rinde la cara trasera respecto de la frontal & adimensional \\
$P$ & Potencia continua medida a la entrada del inversor & W \\
$P_0$ & Potencia nominal instalada del arreglo & Wp \\
$h$ & Horizonte: cuánto hacia adelante se pronostica & s \\
$\mathrm{VI}$ & Índice de variabilidad: cuánto se agita la curva del día & adimensional \\
\bottomrule
\end{tabular}
\end{center}

Nota de notación: el asterisco de $\mathrm{kt^*}$ distingue el índice de \emph{cielo
despejado} del índice de claridad clásico $\mathrm{kt}$, que se define contra la
irradiancia extraterrestre y no contra la de cielo despejado.

% =============================================================================
\section{Radiación}

\subsection{Cielo despejado}

Es la referencia contra la que se mide todo lo demás: cuánta luz habría a esa hora, en
ese punto del planeta, si no hubiera una sola nube. Se calcula con el modelo de Ineichen
con turbidez de Linke climatológica, tal como lo implementa \cd{pvlib}.

\begin{equation*}
  \mathrm{GHI_{cs}} = f_{\text{Ineichen}}(t,\ \text{lat},\ \text{lon},\ \text{alt})
\end{equation*}

\begin{clave}[title=Por qué esto no es hacer trampa]
$\mathrm{GHI_{cs}}$ depende sólo del tiempo y de la posición, no de ningún dato medido.
Por eso se puede calcular para un instante \emph{futuro} sin que eso sea mirar el futuro:
es astronomía, no información privilegiada. Es lo que permite que el pronóstico incorpore
la geometría solar del momento que se está prediciendo.
\end{clave}

\emph{Implementado en:} \cd{predictivo/physics.py} y \cd{agrovoltaic/clearsky.py}.

\subsection{Índice de claridad}

\begin{equation*}
  \mathrm{kt^*} = \frac{\mathrm{GHI}}{\mathrm{GHI_{cs}}}
  \qquad\text{sólo cuando}\quad \mathrm{GHI_{cs}} > 20\ \mathrm{W/m^2}
  \qquad \mathrm{kt^*} \geq 0
\end{equation*}

El umbral evita dividir por casi cero al amanecer y al anochecer, donde el cociente se
dispara sin significar nada. El recorte inferior evita que una medida negativa por
desajuste del sensor produzca un índice negativo, que no existe.

\emph{Implementado en:} \cd{predictivo/physics.py} y la vista \cd{v\_sc\_radiacion\_calibrada}.

\subsection{Corrección del crudo y control de calidad}

El dato crudo se guarda tal cual y se corrige en la capa de consulta, nunca al insertar.
Tres reglas, en este orden:

\begin{equation*}
  \mathrm{GHI_{corr}} =
  \begin{cases}
    0 & \text{si } |\mathrm{GHI_{cruda}} - (-38{,}845)| < 10^{-3} \quad\text{(desajuste nocturno)}\\
    0 & \text{si } \mathrm{GHI_{cruda}} < 0 \\
    \text{NULL} & \text{si } t < \text{2025-07-01} \quad\text{(error de sensor corregido a mediados de 2025)}\\
    \mathrm{GHI_{cruda}} & \text{en caso contrario}
  \end{cases}
\end{equation*}

Y una lectura se marca como sospechosa cuando supera lo que el cielo puede dar:

\begin{equation*}
  \mathrm{qc\_ok} = \mathrm{GHI_{corr}} \leq \max\left(1{,}3 \cdot \mathrm{GHI_{cs}},\ 50\right)
\end{equation*}

El $1{,}3$ deja pasar el realce por reflexión en nubes, que es real y puede superar el
cielo despejado; el piso de 50 W/m$^2$ evita marcar ruido inofensivo de madrugada.

\emph{Implementado en:} \cd{agrovoltaic/ddl.py}, vistas \cd{v\_sc\_radiacion\_corregida} y
\cd{v\_sc\_radiacion\_calibrada}.

% =============================================================================
\section{Pronóstico}

\subsection{La descomposición de la que sale todo}

\begin{equation*}
  \mathrm{GHI} = \mathrm{kt^*} \times \mathrm{GHI_{cs}}
\end{equation*}

De los dos factores, el segundo se calcula con astronomía y el primero es el único que
hay que pronosticar. Pronosticar $\mathrm{kt^*}$ es mucho más fácil que pronosticar
$\mathrm{GHI}$ directo, porque a $\mathrm{kt^*}$ ya se le quitó la curva diurna del sol,
que es la que domina el número.

\subsection{Qué índice se persiste}

No se persiste el $\mathrm{kt^*}$ reciente tal cual, sino su \emph{anomalía} respecto de
lo normal a esa hora, trasplantada a lo normal de la hora que se pronostica:

\begin{equation*}
  \mathrm{kt^*_{pron}} = \mathrm{kt^*_{clima}}(t+h)
  + w \cdot \Big( \mathrm{kt^*_{reciente}} - \mathrm{kt^*_{clima}}(t) \Big)
\end{equation*}

Los dos términos climáticos corrigen cosas distintas y por eso van por separado:

\begin{itemize}[leftmargin=*,itemsep=3pt]
  \item \textbf{El ajuste diurno} (los dos $\mathrm{kt^*_{clima}}$) evita trasladar la
    mañana a la tarde como si fueran iguales. En San Carlos no lo son: el índice medio
    baja de $0{,}58$ a las 11 h a $0{,}35$ a las 16 h, porque las tardes se cierran por
    convección. Ignorarlo producía un sesgo medido de $+52$\,\% a las 16 h.
  \item \textbf{La contracción} (el peso $w$) evita creerle a lo reciente más de lo que
    vale. Persistir supone que la nubosidad de ahora sigue valiendo dentro de $h$, y eso
    se desgasta: la correlación del índice consigo mismo cae de $0{,}90$ a 5 min a
    $0{,}13$ a 6 h.
\end{itemize}

\subsection{El peso}

\begin{equation*}
  w = \rho_h + (1 - \rho_h)\,(1 - r_h)
  \qquad
  r_h =
  \begin{cases}
    0 & h \leq 1\ \mathrm{h} \\
    \dfrac{h - 1\,\mathrm{h}}{3\,\mathrm{h} - 1\,\mathrm{h}} & 1\ \mathrm{h} < h < 3\ \mathrm{h} \\
    1 & h \geq 3\ \mathrm{h}
  \end{cases}
\end{equation*}

donde $\rho_h$ es la autocorrelación del $\mathrm{kt^*}$ medida al desfase $h$ sobre el
histórico anterior al instante que se pronostica, acotada a $[0,1]$. El peso \textbf{se
mide, no se declara}. En horizontes cortos $w=1$ y queda persistencia pura; en horizontes
largos $w=\rho_h$ y manda la correlación real; en el medio se interpola. Sin historia
suficiente (menos de 200 pares) el peso vuelve a 1.

\emph{Implementado en:} \cd{predictivo/forecasters/climatologia.py}.

\subsection{La reconstrucción}

\begin{equation*}
  \mathrm{GHI_{pron}}(t+h) = \mathrm{kt^*_{pron}} \times \mathrm{GHI_{cs}}(t+h)
\end{equation*}

Acá entra la geometría solar del instante pronosticado. Es la diferencia con repetir la
última lectura: si $t+h$ cae más cerca del mediodía, el pronóstico sube aunque las nubes
no se muevan.

\subsection{La banda}

No es una desviación típica sino los \textbf{cuantiles 16 y 84 del error real} que el
método cometió a ese mismo horizonte, medidos sobre el histórico:

\begin{equation*}
  \big[\,\mathrm{GHI_{bajo}},\ \mathrm{GHI_{alto}}\,\big] =
  \Big[\, q_{16}\big(\varepsilon_h\big),\ q_{84}\big(\varepsilon_h\big) \,\Big]
  \qquad \varepsilon_h = \mathrm{GHI_{pron}} - \mathrm{GHI_{real}}
\end{equation*}

Se eligieron 16 y 84 porque en una distribución normal encierran exactamente $\pm 1$
desviación típica, y así la banda se lee igual que siempre aunque el error no sea normal.
La banda anterior, basada en la dispersión de la última hora, salía casi del mismo ancho
para todos los horizontes ($158$ a $166$ W/m$^2$) mientras el error real iba de $169$ a
$237$: prometía una precisión que no tenía.

\emph{Implementado en:} \cd{predictivo/forecasters/uncertainty.py}.

\subsection{Humedad de suelo}

No hay curva del sol que descontar, así que no hay nada que descomponer:

\begin{equation*}
  \hat{y}(t+h) = \mathrm{mediana}\big(y_{\text{reciente}}\big)
  \qquad
  \text{banda} = \big[\,q_{16},\ q_{84}\,\big]\ \text{de la misma ventana}
\end{equation*}

Alcanza porque el suelo cambia despacio y está muy correlacionado consigo mismo.

% =============================================================================
\section{Fotovoltaico}

\subsection{De la horizontal al plano del arreglo}

El piranómetro mide en horizontal, pero los paneles no están en horizontal. Hay que
transponer, y para eso primero hay que separar la luz en sus dos componentes:

\begin{equation*}
  (\mathrm{DNI},\ \mathrm{DHI}) = f_{\text{Erbs}}\big(\mathrm{GHI},\ \theta_z,\ n_{\text{día}}\big)
\end{equation*}

\begin{equation*}
  \mathrm{POA} = f_{\text{isotrópica}}\big(\beta,\ \gamma,\ \theta_z,\ \gamma_s,\
  \mathrm{DNI},\ \mathrm{DHI},\ \mathrm{GHI},\ \rho\big)
\end{equation*}

con $\beta$ la inclinación del plano, $\gamma$ su azimut, $\theta_z$ el ángulo cenital
solar y $\gamma_s$ el azimut solar. El albedo $\rho$ es el \textbf{medido en sitio}, no un
valor de tabla; si cae fuera de $[0{,}05,\ 0{,}9]$ se sustituye por $0{,}2$.

\subsection{Bifacialidad}

Los paneles producen por las dos caras, así que la POA efectiva suma la trasera con un
descuento:

\begin{equation*}
  \mathrm{POA_{ef}} = \mathrm{POA}(\beta,\ \gamma)
  + \varphi \cdot \mathrm{POA}(180^\circ - \beta,\ \gamma + 180^\circ)
  \qquad \varphi = 0{,}80
\end{equation*}

La cara trasera se modela como un segundo plano, el opuesto exacto del frontal.

\emph{Implementado en:} \cd{agrovoltaic/performance.py}.

\subsection{Performance Ratio}

Compara lo que el arreglo produjo con lo que cabría esperar de la luz que efectivamente
recibió:

\begin{equation*}
  \mathrm{PR} = \frac{P / P_0}{\mathrm{POA_{ef}} / 1000}
  \qquad
  P_0 = 1420\ \mathrm{Wp\ por\ arreglo}
  \qquad
  \text{sólo si } \mathrm{POA_{ef}} > 100\ \mathrm{W/m^2}
\end{equation*}

El $1000$ es la irradiancia de las condiciones de ensayo estándar, que es contra lo que
está definida la potencia nominal. El umbral de 100 W/m$^2$ evita calcular el cociente de
noche y con luz muy baja, donde se vuelve inestable.

\emph{Implementado en:} vista \cd{v\_sc\_performance}, generada por \cd{agrovoltaic/ddl.py}.

\subsection{Energía}

Integración trapezoidal de la potencia, o de la irradiancia para la insolación:

\begin{equation*}
  E = \sum_i \frac{P_i + P_{i+1}}{2}\,\Delta t_i
  \qquad [\mathrm{Wh}]
\end{equation*}

% =============================================================================
\section{Calidad de datos}

\subsection{Las dos completitudes}

Son dos preguntas distintas y por eso son dos números. La primera: ¿el registrador estuvo
encendido todas las horas con sol?

\begin{equation*}
  \text{cobertura} = \frac{\text{horas con dato}}{\text{horas de sol del día}}
\end{equation*}

Medir contra 24 horas daría un 50\,\% permanente que no informa nada, porque de noche no
hay nada que medir. Las horas de sol salen del amanecer y el atardecer reales de esa
fecha en ese sitio.

La segunda: dentro de lo que sí grabó, ¿están todas las muestras?

\begin{equation*}
  \text{densidad} = \frac{\text{filas presentes}}{\dfrac{\text{horas con dato} \times 3600}{\text{cadencia}} + 1}
\end{equation*}

La cadencia se infiere del propio día, tomando el intervalo más frecuente, y no de una
constante: el sistema cambió de cadencia varias veces a lo largo de su vida. Un día puede
tener cobertura $1{,}0$ y densidad $0{,}4$: grabó de punta a punta, pero con huecos.

\emph{Implementado en:} \cd{historico/calidad/barrido.py} y \cd{historico/calidad/sol.py}.

\subsection{Índice de variabilidad}

El índice de claridad dice \emph{cuánta} luz llegó, pero no \emph{cómo}. Un día de
$\mathrm{kt^*}$ medio $0{,}5$ puede ser una capa uniforme de nubes o un sol entrando y
saliendo cada dos minutos, y para un inversor no son lo mismo. Eso lo separa el largo de
arco de la curva, comparado con el del cielo despejado:

\begin{equation*}
  \mathrm{VI} = \frac{\sum_i \sqrt{(\Delta \mathrm{GHI}_i)^2 + (\Delta t_i)^2}}
                     {\sum_i \sqrt{(\Delta \mathrm{GHI_{cs}}_i)^2 + (\Delta t_i)^2}}
\end{equation*}

Un día perfectamente despejado da $\mathrm{VI} \approx 1$; el paso de nubes lo dispara.

\begin{aviso}[title={El índice mide la curva, no el registrador}]
El $\Delta t$ tiene que ser el real entre muestras. Con cadencias distintas el mismo día
da índices distintos: se midieron $4{,}15$ y $23{,}81$ para el mismo $\mathrm{kt^*}$ sólo
por cambiar la frecuencia de muestreo. Comparar índices de días con cadencias distintas
no significa nada.
\end{aviso}

\subsection{Clasificación del día}

\begin{center}
\small
\begin{tabular}{@{}p{46mm}p{40mm}p{50mm}@{}}
\toprule
\textbf{Condición} & \textbf{Clase} & \textbf{Qué significa} \\
\midrule
$\mathrm{kt^*}$ alto y $\mathrm{VI}$ bajo & despejado & Sol limpio todo el día \\
$\mathrm{kt^*}$ bajo y $\mathrm{VI}$ bajo & cubierto & Capa uniforme de nubes \\
$\mathrm{VI}$ alto & variable & Intermitente, el peor caso para el inversor \\
resto & parcial & \\
\bottomrule
\end{tabular}
\end{center}

\subsection{Lo físicamente imposible}

\begin{equation*}
  \mathrm{kt^*} > 1{,}2 \ \Longrightarrow\ \text{descartado}
\end{equation*}

Significaría que llegó más energía de la que manda el sol con el cielo abierto. Sin este
filtro, los días con lecturas imposibles se disfrazaban de días soleados: uno llegó a
reportar un índice medio de $5{,}67$.

\subsection{Valores atípicos}

Puntuación robusta, con mediana y desviación absoluta mediana en lugar de media y
desviación típica, para que los propios atípicos no inflen el umbral que debería
detectarlos:

\begin{equation*}
  z = \frac{0{,}6745\,\big(x - \mathrm{mediana}(x)\big)}{\mathrm{MAD}(x)}
  \qquad
  |z| > 3{,}5 \ \Longrightarrow\ \text{atípico}
\end{equation*}

El $0{,}6745$ es el cuantil $0{,}75$ de la normal estándar, y está para que la puntuación
se lea en la misma escala que un $z$ clásico.

\emph{Implementado en:} \cd{predictivo/anomalias.py}.

% =============================================================================
\section{Evaluación}

Todas se calculan sobre el histórico, simulando el reloj: el pronosticador nunca ve datos
posteriores al instante que se le pide predecir.

\begin{equation*}
  \mathrm{MAE} = \frac{1}{n}\sum_i \big| \hat{y}_i - y_i \big|
  \qquad
  \text{sesgo} = \frac{1}{n}\sum_i \big( \hat{y}_i - y_i \big)
  \qquad
  \text{error relativo} = \frac{\mathrm{MAE}}{\bar{y}}
\end{equation*}

El sesgo va aparte del error absoluto a propósito: un método puede tener poco error medio
y aun así equivocarse siempre hacia el mismo lado, que es un defecto distinto y se corrige
distinto.

\begin{equation*}
  \text{destreza}\ [\%] = \left(1 - \frac{\mathrm{MAE}}{\mathrm{MAE_{ingenua}}}\right) \times 100
\end{equation*}

La referencia es la persistencia ingenua, repetir la última lectura. Una destreza de cero
significa que el método no aporta nada sobre no hacer nada.

\emph{Implementado en:} \cd{predictivo/backtest.py}.

% =============================================================================
\section{Constantes}

\begin{center}
\small
\begin{tabular}{@{}p{40mm}p{30mm}p{68mm}@{}}
\toprule
\textbf{Constante} & \textbf{Valor} & \textbf{De dónde sale} \\
\midrule
Latitud, longitud & $10{,}33$, $-84{,}42$ & San Carlos \\
Altitud & 600 m & \\
Zona horaria & UTC$-6$ & Costa Rica, sin horario de verano \\
$P_0$ por arreglo & 1420 Wp & $4 \times 355$ Wp; 2840 Wp en total \\
PV1: inclinación, azimut & $20^\circ$, $150^\circ$ & Arreglo inclinado \\
PV2: inclinación, azimut & $90^\circ$, $50^\circ$ & Arreglo vertical \\
$\varphi$ & $0{,}80$ & Asumido; lo respalda la convergencia del PR \\
Umbral de cielo despejado & 20 W/m$^2$ & Bajo eso, el índice no significa nada \\
Índice máximo creíble & $1{,}3$ & Por encima, se marca la lectura \\
Índice imposible & $1{,}2$ & Por encima, se descarta el día \\
Umbral de POA para el PR & 100 W/m$^2$ & \\
Desajuste nocturno & $-38{,}845$ & Medido en 205 archivos \\
Error del sensor de temperatura & $85{,}0$ °C & Código de error del DS18B20 desconectado \\
Temperatura válida & 10 a 80 °C & Criterio de Leo Cardinale \\
Muestreo eléctrico & 5 min & \\
Muestreo de radiación & 15 s & Mínimo que permite la plataforma de ingesta \\
\bottomrule
\end{tabular}
\end{center}

\vspace{1em}
\begin{clave}[title=Convergencia como validación]
El factor de bifacialidad no se midió, se asumió. Lo que lo respalda es que dos arreglos
con geometrías muy distintas, uno inclinado a $20^\circ$ y otro vertical a $90^\circ$,
aterrizan casi en el mismo Performance Ratio: $0{,}622$ y $0{,}626$. Si el factor
estuviera mal, la ganancia trasera quedaría mal repartida entre los dos y los números se
separarían.
\end{clave}

\end{document}
